% ============================================================
% EvidenceFirst current-results draft
% Generated from repository artifacts available on 2026-06-08.
% Compile from paper/: pdflatex -> bibtex -> pdflatex -> pdflatex
% ============================================================
\documentclass[runningheads]{llncs}

\usepackage[T1]{fontenc}
\usepackage{amsmath,amssymb}
\usepackage{booktabs}
\usepackage{graphicx}
\usepackage{multirow}
\usepackage{xcolor}
\usepackage{xspace}
\usepackage[hidelinks]{hyperref}
\usepackage{microtype}
\usepackage{algorithm}
\usepackage{algpseudocode}

\newcommand{\sys}{\textsc{EvidenceFirst}\xspace}
\newcommand{\eg}{\emph{e.g.},\xspace}
\newcommand{\ie}{\emph{i.e.},\xspace}
\newcommand{\etal}{\emph{et al.}\xspace}
\newcommand{\calG}{\mathcal{G}}
\newcommand{\calP}{\mathcal{P}}
\newcommand{\calE}{\mathcal{E}}

\begin{document}

\title{\sys: Pre-Generation Structural Evidence-State Diagnostics for Traceable KG-RAG}
\titlerunning{\sys: Structural Evidence-State Diagnostics for KG-RAG}

\author{Anonymous Author(s)}
\authorrunning{Anonymous}
\institute{
  Institution anonymized for review\\
  \email{email@anonymized.com}
}

\maketitle

% ============================================================
\begin{abstract}
% ============================================================

Knowledge graph-enhanced retrieval-augmented generation (KG-RAG) is a natural
fit for multi-hop question answering because graph structure can expose the
entity and relation chain needed to answer a question. Yet most KG-RAG systems
inspect evidence too late: they retrieve a graph neighborhood, ask a language
model to generate an answer, and only then evaluate the final string. When the
retrieved graph is incomplete, the system may still produce a fluent answer but
cannot say whether the failure came from a missing entity, a broken bridge
relation, a disconnected component, or the generator itself.

We present \sys, a training-free KG-RAG pipeline that audits structural evidence
state before generation. The central object of inspection is not the answer
string but the graph-structured candidate evidence chain. \sys builds a
per-question evidence graph, checks structural chain completeness, diagnoses
graph gaps, performs localized repair, and then generates a grounded answer from
complete, repaired, or explicitly marked fallback evidence. The output is
therefore both an answer and an audit trail: chain completeness, missing
entities, disconnected pairs, repair attempts, fallback cases, and runtime
failures.

This draft reports the completed comparison under a strict artifact boundary. On
HotpotQA-1000, final \sys v6 reaches 56.2\% EM and 65.6\% F1, descriptively
above HopRAG strict on the full set with the strongest observed gain on bridge
questions. On 2WikiMultiHopQA-500, \sys v6 reaches 69.0\% EM and 75.8\% F1,
obtaining the highest EM among completed runs and outperforming HopRAG strict,
while LightRAG remains slightly higher in F1 and MS GraphRAG for 2Wiki is
deferred.

\keywords{KG-RAG \and GraphRAG \and Evidence Diagnostics \and Multi-Hop QA
\and HotpotQA \and 2WikiMultiHopQA \and Traceable Reasoning}
\end{abstract}

% ============================================================
\section{Introduction}
\label{sec:intro}
% ============================================================

Multi-hop question answering over Web knowledge requires more than retrieving a
passage that shares surface words with the question. A correct answer often
depends on a chain of intermediate entities and relations: an entity mentioned
in the question leads to a bridge page, the bridge page identifies another
entity, and only then can the answer be derived. Retrieval-augmented generation
(RAG) grounds generation in retrieved text~\cite{lewis2020rag}; KG-RAG and
GraphRAG systems add structured entity-relation evidence over the retrieved
corpus~\cite{edge2024graphrag,guo2024lightrag,han2025graphragsurvey}. Graphs
are a natural fit for multi-hop reasoning because the answer is often the
endpoint of a path through evidence.

The open problem is not only retrieval quality. It is evidence sufficiency. A
KG-RAG system may retrieve several relevant triples and still miss the bridge
edge that makes the answer derivable. It may retrieve both endpoint entities
but leave them in disconnected components. It may retrieve a graph neighborhood
whose entities look plausible but do not form the required reasoning chain. In
all of these cases, a retrieve-then-generate pipeline can still produce a
confident answer. The final answer string may be evaluated later, but the
system has already lost the ability to explain whether the failure was caused
by retrieval, graph construction, evidence connectivity, or generation.

\sys addresses this problem by making the evidence graph an inspected
intermediate object. Before generation, the system checks whether the retrieved
graph contains a structurally complete candidate multi-hop chain for the
question. If the chain is incomplete, \sys emits a graph-structural diagnostic:
which entity is missing, which endpoints are disconnected, whether a bridge edge
is absent, or whether evidence retrieval failed entirely. The diagnostic then
drives targeted repair, after which the generator receives complete, repaired,
or explicitly marked fallback evidence.

\paragraph{Contributions.}
This paper makes three contributions.
\begin{itemize}
    \item We introduce pre-generation structural evidence-state diagnostics for
    KG-RAG: the inspected object is the retrieved candidate evidence path, not
    only the generated answer.
    \item We define graph-structural diagnostics for multi-hop QA, including
    chain completeness, missing entities, disconnected pairs, bridge gaps, and
    fallback cases.
    \item We evaluate \sys on HotpotQA-1000 and 2WikiMultiHopQA-500 against
    Naive RAG, IRCoT, A-RAG, LightRAG, MS GraphRAG, and HopRAG where completed,
    and report controlled ablations that separate answer-quality gains from
	    diagnostic traceability.
\end{itemize}

% ============================================================
\section{Related Work}
\label{sec:related}
% ============================================================

\paragraph{RAG and GraphRAG for multi-hop QA.}
RAG conditions generation on retrieved evidence rather than only on parametric
memory~\cite{lewis2020rag}. Dense, active, and self-reflective retrieval methods
improve evidence acquisition~\cite{karpukhin2020dpr,jiang2023active,asai2023selfrag},
while GraphRAG systems organize retrieved information into entity, relation, or
community structures~\cite{edge2024graphrag,guo2024lightrag,han2025graphragsurvey}.
IRCoT interleaves retrieval and reasoning~\cite{trivedi2022ircot}; A-RAG uses an
agentic retrieval trajectory~\cite{arag2026}; and HopRAG performs logic-aware
graph retrieval and pruning~\cite{yu2025hoprag}. \sys is complementary to these
retrieval strategies: it records whether the question-level evidence graph has a
structurally adequate candidate chain before the reader produces the final
answer.

\paragraph{Sufficiency judging and diagnostic evaluation.}
Self-verifying and sufficiency-aware RAG systems use model feedback or external
criteria to reduce unsupported generation~\cite{asai2023selfrag,fairrag2025}.
S2G-RAG judges whether evidence memory is sufficient and produces structured
text gap items~\cite{s2grag2026}; RAGChecker provides fine-grained diagnostic
metrics for retrieval and generation modules~\cite{ru2024ragchecker}. These
systems motivate the same reliability direction but inspect a different object.
\sys narrows the claim to graph-state diagnostics: chain completeness,
missing-entity, disconnected-pair, short-chain, and fallback states are saved
before generation and can condition repair or later audit. These labels are
structural states, not semantic proof that the gold reasoning chain has been
recovered.

% ============================================================
\section{Method}
\label{sec:method}
% ============================================================

\subsection{Pipeline Overview}

Let $q$ be a multi-hop question, $\calP=\{p_i\}$ a set of candidate passages,
and $\calG=(V,E)$ a per-question evidence graph extracted from the passages.
\sys uses six stages:
\begin{enumerate}
    \item build or load a per-question KG from candidate passages;
    \item retrieve KG triples and top-ranked passages for the question;
    \item construct an evidence graph from retrieved triples;
    \item check graph-structural chain completeness;
    \item repair localized graph gaps when possible;
    \item generate a concise answer from verified or repaired evidence.
\end{enumerate}

Algorithm~\ref{alg:evidencefirst} summarizes inference.

\begin{algorithm}[t]
\caption{\sys inference}
\label{alg:evidencefirst}
\begin{algorithmic}[1]
\Require Question $q$, passages $\calP$, per-question KG $\calG$
\Ensure Answer $a$, diagnostic record $D$
\State $D \gets \emptyset$
\State $T \gets$ retrieve\_triples$(q,\calG)$
\If{$T = \emptyset$}
    \State $D \gets D \cup \{\textsc{EmptyEvidence}\}$
    \State $P_q \gets$ local\_passages$(q,\calP)$
    \State $a \gets$ passage\_read$(q,P_q,D)$
\Else
    \State $G_E \gets$ build\_evidence\_graph$(T)$
    \State $(complete, D) \gets$ check\_chain\_completeness$(q,G_E)$
    \If{not $complete$}
        \State $(T', D_r) \gets$ repair\_gap$(q,\calG,D)$
        \State $T \gets T \cup T'$
        \State $G_E \gets$ build\_evidence\_graph$(T)$
        \State $(complete, D') \gets$ check\_chain\_completeness$(q,G_E)$
        \State $D \gets D \cup D_r \cup D'$
    \EndIf
    \State $P_q \gets$ local\_passages$(q,\calP)$
    \State $a \gets$ grounded\_generate$(q,T,P_q,D)$
\EndIf
\State $a \gets$ answer\_canonicalize$(q,a,\calP)$
\State \Return $(a,D)$
\end{algorithmic}
\end{algorithm}

\subsection{Evidence-Chain Diagnostics}

The verifier returns structured diagnostics rather than a single score. Table
~\ref{tab:gap_types} lists the main diagnostic categories. These labels serve
two roles: they guide localized repair during inference, and they provide an
audit trail for posthoc error analysis.

\begin{table}[t]
\centering
\caption{Graph-structural diagnostics emitted by \sys before generation.}
\label{tab:gap_types}
\small
\begin{tabular}{lll}
\toprule
Gap type & Graph condition & Typical repair action \\
\midrule
Missing entity & Required node absent & Entity expansion / passage recall \\
Disconnected pair & Endpoints in separate components & Bridge search \\
Bridge gap & Partial path without bridge edge & Targeted relation retrieval \\
Empty evidence & No usable triples retrieved & Passage fallback \\
\bottomrule
\end{tabular}
\end{table}

The completeness check is structural rather than oracle-semantic. A chain is
treated as complete when the evidence graph contains the question entities and
a connected support structure linking them through retrieved triples or
repaired evidence. A complete chain can still be misread by the generator, but
an incomplete chain warns that the system should repair or fallback before
answer generation.

% ============================================================
\section{Experimental Setup}
\label{sec:setup}
% ============================================================

\paragraph{Datasets.}
We report two evaluation settings. HotpotQA-1000 is a 1000-question
distractor-style multi-hop QA set with repaired contexts from the HotpotQA
validation split. 2WikiMultiHopQA-500 is a 500-question sample whose examples
contain natural-language contexts together with Wikidata-derived evidence
triples, entity ids, evidence ids, and answer ids. Each 2Wiki example includes
10 context passages; the 500 examples expand to 5000 question-context pairs and
3346 unique passage bodies for graph-based baselines.

\paragraph{Baselines.}
For HotpotQA-1000, completed baselines are Naive RAG, IRCoT, A-RAG, LightRAG,
MS GraphRAG, and HopRAG strict response-only rescoring.
For 2Wiki-500, completed baselines are Naive RAG, IRCoT, A-RAG, LightRAG, and
HopRAG strict. MS GraphRAG for 2Wiki was intentionally not run after the user
decision to defer it.

\paragraph{Metrics.}
We use exact match (EM) and token F1 with standard QA normalization:
lowercasing, article removal, punctuation removal, and whitespace collapse. We
report full-sample results only in the main tables. Partial smoke results are
excluded from main comparisons.

\paragraph{Reporting boundary.}
For HotpotQA-1000, we report only the final \sys v6 main model. Earlier
HotpotQA development variants are treated as process artifacts and are omitted
from the main comparison. Controlled ablations use explicit pipeline switches
and are reported separately from historical development variants.

\paragraph{Context-budget and statistical boundary.}
The tables compare completed saved artifacts under their run protocols, not a
context-budget-isolated retrieval benchmark. Reported \sys artifacts can use
local/full reader context, and future strict reruns should use the repository's
\texttt{--no-global-context} switch when the goal is retrieval-budget isolation.
Full-set HotpotQA gains over HopRAG and 2Wiki gains over LightRAG/A-RAG are
descriptive point estimates whose paired confidence intervals cross zero; the
strongest supported accuracy evidence is the HotpotQA bridge subset and the
large 2Wiki gap over HopRAG strict.

\paragraph{Artifact mapping.}
All main numbers in this draft are copied from repository artifacts. The
HotpotQA baseline summary is
\texttt{results/wise\_hotpot1000\_baseline\_summary.csv}. The HotpotQA \sys v6
row uses
\texttt{experiments/evidence\_first/results/hotpotqa\_evidence\_first\_hfctx\_n1000\_splitkg\_v6\_offline\_selector\_guard.csv}.
The 2Wiki baseline rows use
\texttt{results/wise\_2wiki\_naive\_rag\_metrics.csv},
\texttt{results/wise\_2wiki\_ircot\_metrics.csv},
\texttt{results/wise\_2wiki\_arag\_full\_post\_metrics.csv},
\texttt{results/wise\_2wiki\_lightrag\_metrics.csv}, and
\texttt{results/wise\_2wiki\_hoprag\_strict\_metrics.csv}. The 2Wiki \sys row uses
\texttt{results/wise\_2wiki\_evidencefirst\_v6\_readerfull\_canon\_metrics.csv}.
The controlled ablation rows use the 2026-06-08 artifacts with prefix
\texttt{experiments/evidence\_first/results/*evidencefirst\_v6\_ablation\_20260608\_004220\_*}.

% ============================================================
\section{Results}
\label{sec:results}
% ============================================================

\subsection{HotpotQA-1000}

Table~\ref{tab:hotpot_main} reports the completed HotpotQA-1000 baseline
comparison together with the final \sys v6 main model. \sys v6 reaches
56.2\% EM and 65.6\% F1, descriptively 2.1 EM and 1.26 F1 points above HopRAG
strict. The paired full-set confidence interval crosses zero, so this should be
read as competitive completed-run accuracy rather than statistically secure
full-set superiority. Other HotpotQA development attempts are excluded from this
main comparison.

\begin{table}[t]
\centering
\caption{HotpotQA-1000 results. All rows are complete 1000-question runs.}
\label{tab:hotpot_main}
\small
\begin{tabular}{lccc}
\toprule
Method & N & EM & F1 \\
\midrule
Naive RAG & 1000 & 0.3920 & 0.4865 \\
IRCoT & 1000 & 0.4660 & 0.5707 \\
A-RAG & 1000 & 0.5170 & 0.6273 \\
LightRAG & 1000 & 0.4380 & 0.5376 \\
MS GraphRAG & 1000 & 0.1890 & 0.2774 \\
HopRAG strict & 1000 & 0.5410 & 0.6438 \\
\midrule
\sys v6 & 1000 & \textbf{0.5620} & \textbf{0.6564} \\
\bottomrule
\end{tabular}
\end{table}

The result supports a conservative claim: pre-generation structural
evidence-state diagnostics can remain competitive with strong graph baselines
while adding diagnostic traceability. The HotpotQA claim is intentionally
limited to final v6, rather than intermediate attempts explored during
development.

Table~\ref{tab:hotpot_types} shows where the HotpotQA gain comes from. Against
the strongest completed graph baselines, \sys is strongest on bridge questions:
it is 5.8 EM points above HopRAG strict and 10.6 EM points above A-RAG. On
comparison questions, A-RAG and HopRAG strict remain stronger. This is
consistent with the method design: structural chain-state diagnostics are most
useful when the answer depends on connecting entities across evidence, while
attribute-comparison questions can often be solved by direct answer comparison
once the relevant passages are retrieved.

\begin{table}[t]
\centering
\caption{HotpotQA-1000 question-type breakdown against the strongest completed
graph baselines.}
\label{tab:hotpot_types}
\small
\begin{tabular}{lcrrrrrr}
\toprule
Type & N & \multicolumn{2}{c}{A-RAG} & \multicolumn{2}{c}{HopRAG strict} & \multicolumn{2}{c}{\sys v6} \\
\cmidrule(lr){3-4}\cmidrule(lr){5-6}\cmidrule(lr){7-8}
 & & EM & F1 & EM & F1 & EM & F1 \\
\midrule
Bridge & 500 & 0.3500 & 0.4893 & 0.3980 & 0.5221 & \textbf{0.4560} & \textbf{0.5689} \\
Comparison & 500 & 0.6840 & 0.7652 & \textbf{0.6840} & \textbf{0.7654} & 0.6680 & 0.7439 \\
\bottomrule
\end{tabular}
\end{table}

\subsection{2WikiMultiHopQA-500}

Table~\ref{tab:2wiki_main} reports completed 2Wiki-500 results. \sys v6
reader-full canonicalization obtains the highest EM among completed systems,
with point estimates 12.6 EM points above HopRAG strict, 1.6 EM points above
LightRAG, 2.4 EM points above A-RAG, and 3.4 EM points above IRCoT. LightRAG
remains slightly stronger in F1, exceeding \sys by 1.35 F1 points, while \sys is
14.24 F1 points above HopRAG strict. The LightRAG/A-RAG paired intervals cross
zero, so the safe claim is highest completed EM point estimate plus a large
completed-run advantage over HopRAG strict, not universal graph-baseline
dominance.

\begin{table}[t]
\centering
\caption{2WikiMultiHopQA-500 results. All rows are complete 500-question runs.}
\label{tab:2wiki_main}
\small
\begin{tabular}{lccc}
\toprule
Method & N & EM & F1 \\
\midrule
Naive RAG & 500 & 0.2980 & 0.3338 \\
IRCoT & 500 & 0.6560 & 0.7301 \\
A-RAG & 500 & 0.6660 & 0.7323 \\
HopRAG strict & 500 & 0.5640 & 0.6157 \\
LightRAG & 500 & 0.6740 & \textbf{0.7716} \\
\midrule
\sys v6 & 500 & \textbf{0.6900} & 0.7581 \\
\bottomrule
\end{tabular}
\end{table}

The 2Wiki setting is useful for auditing because examples include gold evidence
triples and entity ids. The answer-level result should be read together with the
offline evidence and proxy-adjudication analyses below: they show that
structural graph states are inspectable but only partially aligned with semantic
gold-chain proxies.

\subsection{Component Ablations}
\label{sec:ablations}

Table~\ref{tab:ablations} reports controlled component ablations on both
datasets. The results separate components that directly drive answer quality
from components that primarily provide traceability and failure diagnosis.
Removing local/full reader context causes the largest drop on both datasets,
showing that graph evidence alone is not sufficient for robust answer surface
realization. Removing answer refinement hurts HotpotQA but increases 2Wiki F1
while preserving 2Wiki EM, suggesting that canonical answer selection can
improve exact normalization while reducing partial token overlap. Chain checking
and repair show mixed answer-level effects: repair improves EM relative to the
no-repair setting, but the chain-check switch does not monotonically improve
EM/F1 in the current implementation. We therefore interpret these ablations as
behavioral evidence, not as a claim that every diagnostic stage independently
increases final-answer accuracy.

\begin{table}[t]
\centering
\caption{Controlled \sys ablations. Full rows use final v6; ablations use the
2026-06-08 explicit ablation switches.}
\label{tab:ablations}
\small
\begin{tabular}{lrrrr}
\toprule
\multirow{2}{*}{Variant} & \multicolumn{2}{c}{HotpotQA-1000} & \multicolumn{2}{c}{2Wiki-500} \\
\cmidrule(lr){2-3}\cmidrule(lr){4-5}
 & EM & F1 & EM & F1 \\
\midrule
Full \sys v6 & 0.5620 & 0.6564 & 0.6900 & 0.7581 \\
w/o chain check + repair & 0.5680 & 0.6647 & 0.6780 & 0.7579 \\
w/o graph repair & 0.5540 & 0.6580 & 0.6500 & 0.7424 \\
w/o full/local reader context & 0.2360 & 0.2899 & 0.1840 & 0.2289 \\
w/o answer refinement & 0.5510 & 0.6485 & 0.6900 & 0.7782 \\
\bottomrule
\end{tabular}
\end{table}

\subsection{Selector Calibration}

Table~\ref{tab:selector_calibration} analyzes the final answer-selection signal
inside \sys v6. The selector is not a standalone accuracy guarantee, and the
ablation in Table~\ref{tab:ablations} shows that answer refinement interacts
with dataset-specific EM/F1 normalization. However, the selected cases are much
more accurate than non-selected cases on both datasets. This supports a narrower
claim: answer selection is a useful calibration signal for auditing when the
pipeline trusts the refined evidence-grounded answer.

\begin{table}[t]
\centering
\caption{\sys v6 selector calibration. ``Selected'' means that the final
evidence-grounded postprocess answer was chosen.}
\label{tab:selector_calibration}
\small
\begin{tabular}{lcrrrr}
\toprule
Dataset & Selected? & N & Rate & EM & F1 \\
\midrule
HotpotQA-1000 & Yes & 816 & 0.816 & \textbf{0.6250} & \textbf{0.7346} \\
HotpotQA-1000 & No & 184 & 0.184 & 0.2826 & 0.3096 \\
2Wiki-500 & Yes & 352 & 0.704 & \textbf{0.7358} & \textbf{0.8017} \\
2Wiki-500 & No & 148 & 0.296 & 0.5811 & 0.6546 \\
\bottomrule
\end{tabular}
\end{table}

The 2Wiki canonicalizer is also conservative: it changes only 18 of 500
predictions, and all 18 changed predictions are exact matches. We therefore use
this analysis as supporting evidence for auditability and normalization, not as
a replacement for the controlled answer-refinement ablation.

\subsection{2Wiki Variants}

Table~\ref{tab:2wiki_variants} reports the completed \sys variants on
2Wiki-500. The low scores for the fresh-KG CPU and local-context-only variants
show that answer quality depends strongly on letting the reader see the full
local context. The reader-full canonical variant is the current best 2Wiki
\sys artifact.

\begin{table}[t]
\centering
\caption{Completed \sys variants on 2Wiki-500.}
\label{tab:2wiki_variants}
\small
\begin{tabular}{lccc}
\toprule
Variant & N & EM & F1 \\
\midrule
v6 default & 500 & 0.3920 & 0.4582 \\
v6 fresh KG CPU & 500 & 0.3940 & 0.4565 \\
v6 local context 2 & 500 & 0.5040 & 0.5945 \\
v6 reader-full & 500 & 0.6540 & 0.7470 \\
v6 reader-full canonical & 500 & \textbf{0.6900} & \textbf{0.7581} \\
\bottomrule
\end{tabular}
\end{table}

\subsection{2Wiki Question-Type Analysis}

Table~\ref{tab:2wiki_types} compares \sys v6 and LightRAG by 2Wiki question
type. \sys gains most clearly on pure comparison questions. LightRAG remains
stronger on bridge-comparison F1 and inference F1, suggesting that \sys's
canonical answer selection helps exact matching but can under-preserve partial
answer overlap.

\begin{table}[t]
\centering
\caption{2Wiki-500 question-type breakdown for LightRAG and \sys v6.}
\label{tab:2wiki_types}
\small
\begin{tabular}{lcrrrr}
\toprule
Type & N & \multicolumn{2}{c}{LightRAG} & \multicolumn{2}{c}{\sys v6} \\
\cmidrule(lr){3-4}\cmidrule(lr){5-6}
 & & EM & F1 & EM & F1 \\
\midrule
Bridge comparison & 112 & 0.9464 & 0.9589 & 0.9375 & 0.9375 \\
Comparison & 119 & 0.8992 & 0.9034 & \textbf{0.9580} & \textbf{0.9652} \\
Compositional & 218 & 0.4450 & \textbf{0.6041} & \textbf{0.4541} & 0.5626 \\
Inference & 51 & 0.5294 & \textbf{0.7686} & 0.5294 & 0.7167 \\
\bottomrule
\end{tabular}
\end{table}

\subsection{2Wiki Evidence Availability}

Because 2Wiki provides annotated support facts and evidence triples, we also run
an offline evidence-availability analysis. This analysis does not call an LLM.
It asks whether the evidence visible to the answer reader contains the gold
support titles and evidence entities. Table~\ref{tab:2wiki_support_availability}
compares \sys v6 reader-full against deterministic or logged evidence views
available for Naive RAG and A-RAG. LightRAG is not included in this table because
the completed LightRAG artifacts do not expose comparable per-query retrieved
support titles.

\begin{table}[t]
\centering
\caption{2Wiki-500 evidence availability. A-RAG actual read uses logged
\texttt{chunks\_read\_ids}; A-RAG search upper bound uses all logged keyword-search
hits.}
\label{tab:2wiki_support_availability}
\small
\begin{tabular}{lrr}
\toprule
Evidence view & Support-title recall & Evidence-entity coverage \\
\midrule
\sys v6 reader-full input & \textbf{1.0000} & \textbf{0.9468} \\
Naive BM25 top-5 input & 0.6020 & 0.6685 \\
A-RAG actual read chunks & 0.2955 & 0.3326 \\
A-RAG keyword-search upper bound & 0.5885 & 0.6085 \\
\bottomrule
\end{tabular}
\end{table}

The result supports a narrow evidence-availability claim: \sys v6 does not
filter away annotated 2Wiki support passages before answer generation, whereas
BM25 top-5 and A-RAG's actually read chunks often expose only part of the gold
support set. The same offline analysis also shows diagnostic selectivity:
chain-complete cases reach 70.29\% EM and 76.81\% F1, while disconnected-gap
cases fall to 45.45\% EM and 53.51\% F1 and short-chain cases fall to 63.64\%
EM and 70.35\% F1. We therefore use this table as auxiliary evidence for
traceability and evidence availability, not as a full retrieval leaderboard.

\subsection{2Wiki Gap-Label Proxy Adjudication}

To test whether structural gap labels correspond to semantically meaningful
evidence gaps, we run a no-LLM three-proxy-reviewer adjudication over 100
stratified 2Wiki examples. The proxy views inspect each saved row from strict
gold-chain, reader-support, and conservative meta-review perspectives. This is a
reproducible audit over saved artifacts, not professional human annotation. The
result is deliberately mixed: 40 rows match the structural label, 31 mismatch,
and 29 are ambiguous; the overall match rate is 0.4000 with Wilson 95\% CI
$[0.3094,0.4980]$. Inter-proxy agreement is low, with 72\% of rows showing proxy
disagreement and Fleiss $\kappa=0.1836$. The strongest label is
missing-entities, with a 0.7429 match rate; complete, disconnected, and
short-chain labels are often ambiguous or mismatched. We therefore treat gap
labels as structural audit states whose semantic validity can be inspected, not
as semantic-chain correctness certificates.

\subsection{Deferred Baselines}

Table~\ref{tab:deferred} records the status of runs that are not yet part of
the completed comparison. These entries should not be used for claims until
their full metrics are available.

\begin{table}[t]
\centering
\caption{Deferred 2Wiki-500 baseline.}
\label{tab:deferred}
\small
\begin{tabular}{lll}
\toprule
Method & Status & Notes \\
\midrule
MS GraphRAG & Deferred & User decision: do not run for now. \\
\bottomrule
\end{tabular}
\end{table}

% ============================================================
\section{Discussion}
\label{sec:discussion}
% ============================================================

\paragraph{Accuracy is competitive, but traceability is the main claim.}
On HotpotQA-1000, \sys is descriptively above HopRAG strict on the full saved
artifact and strongest on bridge questions, while adding graph-level diagnostics.
On 2Wiki-500, \sys obtains the highest completed EM point estimate, although
LightRAG remains higher in F1 and MS GraphRAG is still deferred. The robust
claim is therefore not an unqualified leaderboard claim; it is that structural
evidence-state diagnostics can remain competitive while making retrieval
failures inspectable.

\paragraph{Question-type and selector analyses sharpen the claim.}
The HotpotQA gain is concentrated on bridge questions, where \sys has the
largest observed paired advantage and where evidence-chain connection is
directly tested. The selector analysis shows that answers chosen by the final
evidence-grounded selection step are substantially more accurate than
non-selected answers on both datasets. These results support \sys as an
auditable pipeline with useful internal confidence signals, rather than a
method that uniformly dominates every question type.

\paragraph{2Wiki exposes canonicalization pressure.}
The 2Wiki results show a split between EM and F1. \sys v6 has the highest EM
point estimate through reader-full context and canonical answer selection, but
LightRAG achieves higher F1. This suggests that some \sys errors are near misses
or overly canonicalized outputs. The next revision should inspect cases where
\sys loses F1 but not EM, and cases where LightRAG has high F1 without exact
match.

\paragraph{Ablations support a narrower component claim.}
The controlled ablations do not support a simple monotonic story in which every
EvidenceFirst component independently improves answer EM/F1. The strongest
accuracy component is local/full reader context. Answer refinement is useful on
HotpotQA but is not uniformly beneficial on 2Wiki F1. Chain checking and repair
remain central to the method's traceability claim because they expose missing
entities, disconnected evidence, bridge gaps, and repair attempts before
generation, but their current answer-level gains are mixed and should be
presented as diagnostic behavior rather than guaranteed accuracy improvement.

\paragraph{HotpotQA uses one final model.}
The HotpotQA section reports final \sys v6 as the single main model. Earlier
HotpotQA development attempts are not part of the reported method comparison
or ablation evidence.

\paragraph{2Wiki graph baselines are only partly complete.}
HopRAG strict is now included in the 2Wiki table and is below \sys on both EM
and F1. MS GraphRAG remains deferred, so claims about graph-baseline dominance
on 2Wiki should still remain cautious.

% ============================================================
\section{Limitations}
\label{sec:limitations}
% ============================================================

This draft has five limitations. First, the main tables are saved-protocol
artifact comparisons, not context-budget-isolated retrieval benchmarks; future
strict reruns should use \texttt{--no-global-context}. Second, full-set paired
intervals cross zero for the smaller HotpotQA and 2Wiki differences, so those
rows should be treated as descriptive point estimates. Third, MS GraphRAG was
intentionally deferred on 2Wiki, so that baseline remains absent. Fourth, the
current ablations are system-level switches, and some components interact:
reader context can compensate for imperfect structural graph checks, while
answer canonicalization changes EM and F1 differently. Fifth, the 100-example
semantic gap-label adjudication is a reproducible proxy audit rather than a
professional human annotation study, and its 72\% disagreement rate shows that
structural labels should not be treated as semantic correctness labels.

% ============================================================
\section{Conclusion}
\label{sec:conclusion}
% ============================================================

We introduced \sys, a KG-RAG pipeline that records structural evidence-chain
state before generation. The method changes the reliability object from only the
final answer string to the retrieved evidence graph, enabling graph-structural
diagnostics such as missing entities, disconnected pairs, bridge gaps, and
fallback cases. Current results show that \sys is competitive on HotpotQA-1000
and obtains the highest completed EM point estimate on 2Wiki-500. More
importantly, it provides a traceable account of where evidence succeeds or
fails, moving KG-RAG toward systems that are not only accurate, but auditable.

\bibliographystyle{splncs04}
\bibliography{comagraag}

\end{document}
